\begin{abstract}
A run-to-run standard error for a benchmark average built from per-benchmark variances ignores covariance between benchmarks, and in simulation a margin-scale comparison using it declares a difference in 0.113 of null replicates, against a nominal 0.05. SNAP estimates it from cross-half products of run-centred item-half scores on a fixed battery of likelihood-scored multiple-choice benchmarks. On 375 released DataDecide runs at the largest checkpoint each configuration's runs share, the ten-benchmark per-byte margin average has 1.244 times its independent standard deviation \ci{1.143}{1.338}. BoolQ carries 68\% of the covariance trace, and single-benchmark removals give 1.096 to 1.786. A uniform item component shared across benchmarks would need more than all released item-noise variance to produce 1.244. The accuracy interval \ci{0.993}{1.157} includes one, but this design detects a true 1.10 in 0.135 of replicates. A pre-specified test on four unseen tasks fails its lower-limit rule at 1.217 \ci{0.872}{1.498}. \outcome{\Ronecase}{1}{R1 pass}{\outcome{\Rtwocase}{1}{R2 not supported}{Under rules fixed before scoring, nine seeds at five PolyPythias sizes give margin inflation of \pending{R1 lambda} with interval \ci{\pending{lo}}{\pending{hi}}, and pairing that battery with a disjoint item bank leaves \pending{cross lambda}, which gives no support to the shared-item explanation.}\outcome{\Rtwocase}{2}{R2 supported}{Under a rule fixed before scoring, nine seeds at five PolyPythias sizes give margin inflation of \pending{R1 lambda} with interval \ci{\pending{lo}}{\pending{hi}}, although a disjoint item bank supports a shared item component.}\outcome{\Rtwocase}{3}{R2 undecided}{Under a rule fixed before scoring, nine seeds at five PolyPythias sizes give margin inflation of \pending{R1 lambda} with interval \ci{\pending{lo}}{\pending{hi}}, although a disjoint item bank leaves the source of that inflation undecided.}}\outcome{\Ronecase}{2}{R1 fail}{On a second family, nine seeds at five PolyPythias sizes give margin inflation of \pending{R1 lambda} with interval \ci{\pending{lo}}{\pending{hi}}, which fails a rule fixed before scoring at a design with 40 within-configuration contrasts, and a disjoint item bank \outcome{\Rtwocase}{1}{R2 not supported}{leaves \pending{cross lambda}, which gives no support to the shared-item explanation}\outcome{\Rtwocase}{2}{R2 supported}{supports a shared item component}\outcome{\Rtwocase}{3}{R2 undecided}{leaves the source of that inflation undecided}.} Beyond the pre-specified test and two PolyPythias rules, all analyses are exploratory. We recommend reporting paired replicate scores with a benchmark-removal check.
\end{abstract}

\section{Introduction}
A benchmark average can change across training runs even when the evaluation items stay fixed. For an equally weighted average of $K$ benchmarks, independent run deviations give variance $K^{-2}\operatorname{tr}\Sig$, where $\Sig$ is their cross-benchmark covariance matrix. Covariance adds the off-diagonal terms to that variance, and assuming independence therefore weakens a comparison's error control. In simulation with margin-like covariance, such a comparison of two independently trained configurations on the margin scale declares a difference in 0.113 of replicates with a true gap of zero, against a nominal rate of 0.05 (Appendix~\ref{app:supplementary}). On DataDecide's own recipe comparisons the correction changes few calls (Section~\ref{sec:prediction}), so the understatement shows mainly in the error bar of one model's battery average and in planning replicate runs.

We estimate item-general covariance from two disjoint item halves per benchmark, centring both half scores across the replicate runs within each configuration and taking their cross-half products. Under zero-mean item noise and conditional independence across halves, they estimate run covariance, including its diagonal. We adapt the repeated-measurement construction of quantitative genetics to this setting and call it SNAP, Seed Noise Across Phenotypes (Appendix~\ref{app:derivations}).

Our empirical study uses DataDecide \citep{magnusson2025}, with 25 data recipes and five of its model sizes, each with three released replicates, and those 125 configurations supply 375 runs across ten benchmarks, scored at the largest checkpoint step each configuration's replicates share. Our primary analysis reads released item scores without model inference and compares per-byte log-likelihood margins with accuracy on the same items. These two scales respond differently to changes in confidence, and margins carry more seed signal, with per-benchmark reliability averaging 0.652 against 0.324 for accuracy (Section~\ref{sec:composition}), which is why we treat margins as the primary scale and read accuracy as a bound.

Margin inflation $\Lam$, the ratio of the average's standard deviation to its independent value, is 1.244 with recipe-cluster interval \ci{1.143}{1.338}, which excludes one, while accuracy gives 1.078 \ci{0.993}{1.157}, which includes one at the scored checkpoint (Table~\ref{tab:primary}).

Because we built our estimator on these ten benchmarks, we pre-specified a test on SciQ, MedMCQA, and multiple-choice DROP and CoQA on the same checkpoints at 530M, 750M, and 1B, whose rule requires the lower interval limit for held-out margin inflation to exceed one. The test fails on all 225 runs, where held-out inflation is 1.217 \ci{0.872}{1.498} (Section~\ref{sec:heldout}).

A single model family and a seed label that also changes the training batch leave open whether the inflation belongs to DataDecide's design. We therefore score all 45 PolyPythias runs at five sizes on the same 37,682 items, where each of nine seeds reruns one recipe, and pair that battery with a disjoint bank of 6,808 items built with the OLMES harness at a pinned commit. Two rules fixed before scoring read the results, and in simulation the first passes in only 0.259 of replicates at a true 1.244 and the second returns undecided in 0.761, so we read both as bounds. Section~\ref{sec:family} reports that margin inflation there is \pending{R1 lambda} with interval \ci{\pending{lo}}{\pending{hi}} and that the shared-item explanation is \outcome{\Rtwocase}{1}{R2 not supported}{not supported}\outcome{\Rtwocase}{2}{R2 supported}{supported}\outcome{\Rtwocase}{3}{R2 undecided}{undecided}.

This paper makes four contributions to measuring run covariance. We define the item-general run covariance that cross-half products estimate, identified up to an item component shared across benchmarks and halves. Unlike the battery average's replicate standard deviation, this target transfers to redrawn banks and separates run variance from item noise. Algorithm~\ref{alg:snap} gives the measurement recipe with near-nominal intervals, whose wild recipe-cluster coverage runs from 0.928 to 0.954 across twelve simulated populations. We measure it on two score scales across the 375 released runs, where margin inflation is 1.244 and every recipe, recipe-pair and benchmark deletion, along with four of the five size-band deletions, leaves an interval that excludes one. We also pre-specify a test on unseen tasks and report that it fails at its pre-specified scope.
